% =============================================================================
% Section 9: Open problems
% Crisp, well-scoped, attractive to outside contributors. No overclaiming.
% =============================================================================

\section{Open problems}
\label{sec:open}

The structural reduction and the finite-local barrier of this paper
make precise what remains genuinely open about Erd\H{o}s
Problem~\#647. We collect five problems whose resolution would
either close the residual frontier, refine the conditional picture,
or rule out further finite-local progress. Each is stated as a
self-contained question, with the relevant context and the
sharpest known partial result.

\subsection{The Hensel load-balancing conjecture}
\label{sec:open-hensel-balance}

The kernel framework of \cref{sec:hensel-kernel-audit} produces, for
every $r \in \Ropen$ and every protected window $K_0$, a
Hensel-complete CRT class $a \pmod{\Lcap}$ satisfying Conjecture~B
$(2 \tauf(D_k(a)) \le k + 2)$ at every shift $1 \le k \le K_0$. The
empirical observation through $K_0 \in \{95, 200, 1021\}$ and the
sentinel runs at $K_0 \in \{1333, 2162, 2756, 3481\}$ on the Hensel
compute box (cf.\ \cref{sec:empirical-detail}) is that
admissibility persists indefinitely.

\begin{problem}[Hensel load-balancing]
\label{prob:hensel-balance}
Prove or disprove: for every $r \in \Ropen$ there exists a
Hensel-complete kernel $(a, \Lcap)$ such that Conjecture~B holds for
\emph{all} $k \ge 1$.
\end{problem}

A positive resolution would upgrade the per-window kernel data of
\cref{sec:hensel-kernel-audit} to a structural barrier statement
quantifying over all shifts simultaneously, and would close the
``barrier'' row of \cref{tab:audit-summary} in the negative
direction. A negative resolution would identify the first $k$ at
which Hensel admissibility fails, providing a concrete shift bound
above which the prime-residual kernel framework cannot certify
survivors. Both resolutions have value; the question itself is
quantifiable, finite-checkable for any individual $K_0$, and tightly
coupled to the kernel-space CSP.

\subsection{A non-finite-local closure mechanism}
\label{sec:open-non-local}

\Cref{thm:ple-insufficiency,thm:interval-cover-insufficiency} rule
out closure within the family of finite-local forced-divisor
certificates of \cref{def:finite-local-cert}, at the parameter
regimes tested. \Cref{sec:obstructions} rules out closure within
density-averaged sieve methods of Selberg/Shiu/Nair--Tenenbaum
type, $k \le 2$ correlation methods, and prime-detection methods
of Maynard type. The remaining mathematical territory is large but
unspecified.

\begin{problem}[Non-finite-local closure]
\label{prob:non-local}
Exhibit a closure certificate for at least one $r \in \Ropen$ that
is \emph{not} expressible in the finite-local forced-divisor
language of \cref{def:finite-local-cert} and is also not subsumed
by the analytic obstructions of \cref{sec:obstructions}.
Equivalently: identify a mathematical primitive outside the union
of these two families that yields a closure for some $r \in
\Ropen$.
\end{problem}

The phrasing is deliberately broad. Candidate mechanisms include
Erd\H{o}s covering systems with non-prime-power moduli of a kind not
captured by the kernel CRT decomposition (e.g.\ moduli with shared
prime factors leading to genuinely composite cover obstructions);
deterministic finite covers of $\Ropen$ obtained from a global
arithmetic identity; and non-multiplicative methods that exploit
additive structure of $360360 \cdot Y - k$ rather than its divisor
function. The problem is open even at the level of producing a
single new $r$.

\subsection{Classification of $\Uset$ and the seven missing cells}
\label{sec:open-classify}

The 41-element subset $\Uset \subset \Z / 323\Z$ of
\cref{sec:q323-normalization} has the curious structural property
that its projections to $\Z / 17\Z$ and $\Z / 19\Z$ have sizes 6
and 8 respectively, yet $|\Uset| = 41$ rather than $48 = 6 \cdot 8$.
Of the 48 product cells, exactly 7 are missing, namely $\{188, 233,
275, 278, 281, 284, 290\}$.

\begin{problem}[Classification of $\Uset$]
\label{prob:classify-U}
Identify the structural rule(s) that determine which 41 of the 48
product cells of $\mathrm{proj}_{17}(\Uset) \times
\mathrm{proj}_{19}(\Uset)$ appear in $\Uset$, equivalently, which
7 cells are missing. State the rule purely in terms of the
Stage-1 sieve elimination conditions of
\cref{sec:stage1-sieve}.
\end{problem}

A satisfactory answer would express each missing cell as the joint
solution of a small, named subset of Stage-1 sieve conditions,
ideally drawing on the same structure as the 143-collapse
(\cref{rem:143-collapse-first-principles}). The current Lean
infrastructure can compute the 7 missing cells but does not
identify a structural rule; the problem is well-posed and modest in
scope, and a clean answer would tighten the operational frontier.

\subsection{Pointwise correlation and the CGK obstruction}
\label{sec:open-pointwise-correlation}

The four convergent obstructions of
\cref{sec:synthesis-no-closure} include the $k \le 2$
correlation ceiling: the strongest available $\tauf$-correlation
results of MRT and MSTT operate at $k = 1$ (Tao--Ter\"av\"ainen
2025) and $k = 2$ (MRT 2019, MSTT-I/II 2024) only, in an averaged
$L^1$ sense. The kernel-space CSP requires pointwise estimates at
$k \ge 3$ jointly across an interval, which exceed the reach of
current methods.

\begin{problem}[Pointwise multi-shift $\tauf$-correlation]
\label{prob:pointwise-correlation}
Prove a pointwise $\tauf$-correlation estimate of the form
\[
\tauf(n) \cdot \tauf(n - h_1) \cdots \tauf(n - h_k)
\;\le\; C_{k, \mathbf{h}} \cdot \log^k(n)
\]
\emph{pointwise} for $n$ in a residue class $r \pmod q$ with shift
vector $\mathbf{h} = (h_1, \ldots, h_k)$ and $k \ge 3$, with
$C_{k, \mathbf{h}}$ explicit and small enough to load-bear against
$k + 2$. Alternatively: extend the Conrey--Gonek--Keating
heuristic (CGK) to a quantitative pointwise upper bound at $k \ge
3$ in the residue regime relevant to $\Ropen$.
\end{problem}

The challenge is the gap between $L^1$ averaging
(where MRT/MSTT operate) and pointwise control (where the barrier
needs estimates), compounded by the well-known divergence of
$\tauf$ at smooth $n$. A successful pointwise estimate would
collapse the third obstruction of \cref{sec:synthesis-no-closure}
and reopen the analytic-NT path; a definitive obstruction proof
would close the obstruction picture.

\subsection{The fundamental question: existence of a counterexample}
\label{sec:open-existence}

Despite the structural reduction, the finite-local barrier, and the
empirical floor at the principal residue, the problem of whether
Erd\H{o}s~\#647 admits a counterexample remains open in the strict
sense.

\begin{problem}[Erd\H{o}s~\#647 itself]
\label{prob:e647-itself}
Determine whether there exists $n > 24$ with
\[
\max_{m < n} \bigl( m + \tauf(m) \bigr) \;\le\; n + 2 .
\]
\end{problem}

The conditional analysis of \cref{sec:obstructions} (under
Schinzel's Hypothesis~H) is consistent with an infinite survivor
set with positive Bateman--Horn singular-series density; the
empirical floor of \cref{prop:empirical-r0} establishes no
counterexample for $n \le 3.34 \times 10^{21}$ contiguously at the
principal residue. Neither rules out a counterexample, and neither
constructs one. We regard \cref{prob:e647-itself} as the
fundamental open question.

\subsection{Closing remark}
\label{sec:open-closing}

The five problems above are independent in scope and can be
attacked in any order. \Cref{prob:hensel-balance} is the most
finite-checkable, requires no new analytic input, and is the
natural next milestone within the framework of this paper.
\Cref{prob:classify-U} is modest in scope and accessible to
unaided pen-and-paper analysis. \Cref{prob:non-local} and
\cref{prob:pointwise-correlation} are the deepest open questions
and would require new mathematical input from outside the
finite-local + density-averaged toolkit. \Cref{prob:e647-itself}
is the headline question. The barrier of this paper is intended to
make precise where these problems live, and which of them are
prerequisite to which.

% =============================================================================
% End of section 9
% =============================================================================
